%
% menu_lifecycle_bounded_inbox_nolock_buggy.tex
%
% FIDELITY CONTROL for docs/menu_lifecycle_bounded_inbox.tex.
%
% This spec differs from the fixed model in one place: the inbox lock is
% removed, so BoundedInbox.put is no longer an atomic critical section. Its
% check-then-act -- (1) read qsize, (2) drop the front if full, (3) enqueue --
% splits into three separately-schedulable steps, and the TWO producers (the
% menu-click offer and the publish fan-out) interleave them freely, exactly as
% they do in the shipped code where offer's put runs under _inboxes_lock and the
% writer's put runs with it released.
%
% It reproduces BOTH defects the review found:
%
%   OVERSHOOT (B1's negation, depth > cap): from depth = cap - 1 both producers
%     Check (qsize < cap, neither will drop) then both Enq -- depth = cap + 1.
%       probcli ... -model_check -nodead -goal "depth > cap" => goal FOUND
%       (e.g. OfferCheck; OfferEnqPlain reaching depth 1, then
%        OfferCheck; PubCheck; OfferEnqPlain; PubEnqPlain -> depth 3).
%
%   DOUBLE-DROP (B2's negation, spuriousDrop = fset): from depth = cap both
%     producers Check (qsize >= cap, both will drop); the first DropStep lowers
%     depth to cap - 1, the second DropStep then fires at depth < cap -- a
%     get_nowait discarding a message that had room.
%       probcli ... -model_check -nodead -goal "spuriousDrop = fset" => goal FOUND
%
% Against the fixed spec both goals are NOT found: the lock makes the compound
% atomic, so no producer can interleave another's check, drop, and enqueue.
%
% The two producers are modelled as the two elements of THREAD and the per-thread
% put state as functions THREAD -> PUTPHASE and THREAD -> FSTATE, because without
% the lock the producers are symmetric -- both run the same un-serialized put.
%

\documentclass[a4paper,10pt,fleqn]{article}
\usepackage[margin=1in]{geometry}
\usepackage{fuzz}
\usepackage[colorlinks=true,linkcolor=blue,citecolor=blue,urlcolor=blue]{hyperref}

\begin{document}

\title{lux Bounded Inbox --- No-Lock Fidelity Control}
\author{The check-then-act put torn by two producers without the inbox lock}
\date{September 2026}
\maketitle

This is a fidelity control for \texttt{menu\_lifecycle\_bounded\_inbox.tex}. The
inbox lock is removed and the compound \texttt{put} splits into its three
schedulable steps, run by two symmetric producers.

\begin{axdef}
  cap : \nat \\
  maxDepth : \nat
\where
  cap = 2 \\
  maxDepth = 3
\end{axdef}

\begin{zed}
  FSTATE ::= fclear | fset
\end{zed}

\begin{zed}
  THREAD ::= toffer | tpublish
\end{zed}

\begin{zed}
  PUTPHASE ::= pidle | pchecked | pdropped
\end{zed}

\begin{schema}{Inbox}
  depth : \nat \\
  spuriousDrop : FSTATE \\
  phase : THREAD \fun PUTPHASE \\
  willDrop : THREAD \fun FSTATE
\where
  depth \leq maxDepth
\end{schema}

\begin{schema}{Init}
  Inbox'
\where
  depth' = 0 \\
  spuriousDrop' = fclear \\
  phase' = \{ toffer \mapsto pidle, tpublish \mapsto pidle \} \\
  willDrop' = \{ toffer \mapsto fclear, tpublish \mapsto fclear \}
\end{schema}

% -- Check: read qsize, decide whether the front will be dropped --------------
% Models step (1). The decision is snapshotted into willDrop(t?) and can go
% stale before the drop executes -- the tear the lock forbids.
\begin{schema}{Check}
  \Delta Inbox \\
  t? : THREAD
\where
  phase~t? = pidle \\
  (depth \geq cap \implies willDrop' = willDrop \oplus \{ t? \mapsto fset \}) \\
  (depth < cap \implies willDrop' = willDrop \oplus \{ t? \mapsto fclear \}) \\
  phase' = phase \oplus \{ t? \mapsto pchecked \} \\
  depth' = depth \\
  spuriousDrop' = spuriousDrop
\end{schema}

% -- DropStep: the get_nowait for a thread that decided to drop ----------------
% Models step (2). If depth has fallen below cap since the check, this drop
% discards a message that had room: spuriousDrop is set.
\begin{schema}{DropStep}
  \Delta Inbox \\
  t? : THREAD
\where
  phase~t? = pchecked \\
  willDrop~t? = fset \\
  (depth < cap \implies spuriousDrop' = fset) \\
  (depth \geq cap \implies spuriousDrop' = spuriousDrop) \\
  (depth > 0 \implies depth' = depth - 1) \\
  (depth = 0 \implies depth' = 0) \\
  phase' = phase \oplus \{ t? \mapsto pdropped \} \\
  willDrop' = willDrop
\end{schema}

% -- EnqAfterDrop: the enqueue following a drop -------------------------------
\begin{schema}{EnqAfterDrop}
  \Delta Inbox \\
  t? : THREAD
\where
  phase~t? = pdropped \\
  depth' = depth + 1 \\
  phase' = phase \oplus \{ t? \mapsto pidle \} \\
  spuriousDrop' = spuriousDrop \\
  willDrop' = willDrop
\end{schema}

% -- EnqPlain: the enqueue for a thread that did NOT drop ----------------------
% Models step (3) for a not-full check. Two of these from depth = cap - 1
% overshoot to cap + 1.
\begin{schema}{EnqPlain}
  \Delta Inbox \\
  t? : THREAD
\where
  phase~t? = pchecked \\
  willDrop~t? = fclear \\
  depth' = depth + 1 \\
  phase' = phase \oplus \{ t? \mapsto pidle \} \\
  spuriousDrop' = spuriousDrop \\
  willDrop' = willDrop
\end{schema}

% -- Drain: recv consumes one, so contention recurs ---------------------------
\begin{schema}{Drain}
  \Delta Inbox
\where
  depth > 0 \\
  depth' = depth - 1 \\
  spuriousDrop' = spuriousDrop \\
  phase' = phase \\
  willDrop' = willDrop
\end{schema}

\end{document}
